%=====================================================================
\section{Small dimensions: segment shrinkage and the near-copy guard (v2.2)}\label{sec:g3small}
% numbered from 1720 (gap plan G3 range); the shared counter is restored at the end
\newcounter{savedpropgthree}\setcounter{savedpropgthree}{\value{proposition}}
A small dimension (AdventureWorks \texttt{DimProduct}, 606 rows) is fitted by many latent segments of few rows: its strata give 24 segments of 6 to 49 rows, none large enough for its own copula (\texttt{copula.min\_rows}), so every segment already uses the pooled copula and only its marginals are its own. With the Dirichlet weight $a_0=0.25$ (\texttt{marginals.dirichlet\_prior}) a segment marginal is close to the empirical law of its rows, and synthetic rows land closer to real rows than real rows are to each other (\texttt{dcr\_ratio} 0.40--0.65 at the defaults, limit 0.5). This section gives the two remedies of \texttt{relsynth} 2.2 (guide \texttt{docs/privacy.md}, section ``Small dimensions''; code \texttt{relsynth/generators/copula.py}, \texttt{relsynth/nearcopy.py}). Each proposition is tested by a check in \texttt{tests/science/test\_small\_dims\_checks.py} (C1720--C1724).

\setcounter{proposition}{1719}
\begin{proposition}[Segment shrinkage]\label{prop:segshrink}
Let a segment of an entity table have $n$ rows (for a continuous column, $n$ is the number of its non-null values that the fit keeps, at most \texttt{max\_obs}), empirical law $e$ of a column, and let $F_0$ be the table's marginal of that column. With $a=\max(a_0,\texttt{segment\_prior\_rows})$, the segment marginal is the posterior $\mathrm{DP}(a+n,\,(aF_0+ne)/(a+n))$ for a continuous column and $\mathrm{Dir}(c+ap_0)$ for a discrete one ($c$ the level counts, $p_0$ the table's level probabilities). (a) Its posterior predictive is $w\,e+(1-w)F_0$ with $w=n/(n+a)$: a synthetic value comes from the segment's own rows with probability $w$. (b) The number of random variates drawn does not depend on $a$: one Dirichlet, or one Beta and one Dirichlet. The number of underlying random numbers can: NumPy draws a Beta by rejection, so a different $a$ can leave the generator in a different state, and later draws differ too. Only $a=a_0$ (\texttt{segment\_prior\_rows} at most $a_0$, in particular 0) reproduces the 2.1 stream, and with it the 2.1 output.
\end{proposition}
\begin{proof}
(a) For the discrete case $E[p]=(c+ap_0)/(n+a)=w\,e+(1-w)p_0$. For the continuous case the code draws $V\sim\mathrm{Beta}(n,a)$, Bayesian-bootstrap weights on the $n$ rows and equal weights on the quantile knots of $F_0$ (the representation of the DP posterior in Prop.~\ref{prop:dp}); $E[V]=n/(n+a)=w$. (b) One Dirichlet (discrete) or one Beta plus one Dirichlet over the $n$ rows (continuous) is drawn per segment and column whatever $a$ is. With $a=a_0$ the parameters, and so every draw, are those of 2.1.
\end{proof}
\noindent Check C1720 draws 20{,}000 posterior marginals of a segment for $a\in\{0.25,10,20\}$ and compares the share of predictive mass on the segment's own atoms with $w$, counts the variates drawn by one fit under two values of $a$ (equal), and replays the Beta draw over many seeds to show that the generator's state after it can differ between two values of $a$, which is why (b) speaks of variates and not of the stream.

\setcounter{proposition}{1720}
\begin{proposition}[Large segments are unchanged]\label{prop:seglarge}
In the notation of Prop.~\ref{prop:segshrink}, with the same prior law $F_0$ under both weights, the total variation distance between the posterior predictive laws (the mixtures of (a), before the fitted marginal is put on its grid) under $a$ and under $a_0<a$ is $\frac{n(a-a_0)}{(n+a)(n+a_0)}\,\mathrm{TV}(e,F_0)\le\frac{a-a_0}{n+a}$.
\end{proposition}
\begin{proof}
Both predictives are mixtures $w e+(1-w)F_0$ with weights $w_a=n/(n+a)$ and $w_{a_0}=n/(n+a_0)$, so their difference is $(w_{a_0}-w_a)(e-F_0)$ and its total variation is $|w_{a_0}-w_a|\,\mathrm{TV}(e,F_0)$, with $w_{a_0}-w_a=n(a-a_0)/((n+a)(n+a_0))\le(a-a_0)/(n+a)$ and $\mathrm{TV}\le1$.
\end{proof}
\noindent With the default of 20 rows a segment of 2{,}000 rows moves by at most 1\%, a segment of 30 rows by up to 40\%. Fact tables are not shrunk: their segments are blocks of the relational model, not groups of records. Check C1721 computes both sides for random $e$, $F_0$, $n$ and $a$. When the prior changes with $a$, as in shape mode, Prop.~\ref{prop:segshape}(e) gives the bound $a/(n+a)$ instead.

\setcounter{proposition}{1721}
\begin{proposition}[The near-copy guard]\label{prop:ncguard}
Let $d$ be the metric of \texttt{relsynth.nearcopy} on a table's real rows $X$, $\nu(x)$ the distance from $x\in X$ to its nearest real row at a positive distance, and $B=\{s:\ 0<d(s,x^*)<\rho\,\nu(x^*)\}\cup\{s=x^*:\ x^*\text{ has no twin in }X\}$, $x^*$ the nearest real row to $s$. A row drawn from its segment law $P_g$ and redrawn while it lies in $B$, at most $R$ times, (a) has law $P_g(\cdot\mid B^c)$ when it is accepted, (b) is accepted with probability $1-P_g(B)^{R+1}$, and (c) when no row of the table falls in $B$ the guard draws no random number and the table is unchanged.
\end{proposition}
\begin{proof}
(a) and (b) are rejection sampling with independent proposals from $P_g$: the accepted draw has the conditional law and the number of rejections is geometric. (c) The redraw loop starts only when some row is in $B$.
\end{proof}
\noindent Rows that are still in $B$ after $R$ rounds are kept and counted in the sample log (\texttt{near\_copies}). Check C1722 runs the guard on a fitted table, checks every accepted row against $B$ by brute force, compares the accepted law with $P_g(\cdot\mid B^c)$ on a one-dimensional example, and checks (c) byte for byte.

\setcounter{proposition}{1722}
\begin{proposition}[The guard is not a privacy guarantee]\label{prop:ncnotdp}
Let $M$ be any generator whose law has positive density on a neighbourhood of a record $x$, and $M_B$ the same generator followed by the guard with $R=\infty$. For the neighbouring tables $X$ and $X\cup\{x\}$ (with $x$ not a twin of a row of $X$ and $\nu(x)>0$), $M_B(X)$ puts positive probability on the ball $\{s: d(s,x)<\rho\nu(x)\}$ minus the guard regions of $X$, and $M_B(X\cup\{x\})$ puts none there, unless that set is empty. So $M_B$ is not $\varepsilon$-differentially private for any finite $\varepsilon$, and the absence of synthetic rows around a real record is evidence of the record.
\end{proposition}
\begin{proof}
Under $X\cup\{x\}$ every synthetic row whose nearest real row is $x$ and lies within $\rho\nu(x)$ of it is in $B$ and is never accepted; a row nearer to another real row than to $x$ lies outside that ball. Under $X$ the same set is outside $B$ (the rows of $X$ are farther than their own guard radii, by the choice of the set) and has positive probability by the density assumption. A ratio of probabilities $p/0$ is unbounded.
\end{proof}
\noindent With a finite number of rounds $R$ some near copies are kept, and the proposition does not apply; the guard still adds no guarantee of its own. In DP mode stage~2 is fitted on the DP seed database, so the guard reads seed rows only and is post-processing of the DP release, which keeps its $(\varepsilon,\delta)$. This is the hole that Ganev and De Cristofaro \cite{ganev2023dcr} use to reconstruct records from outputs filtered by distance to closest record. relsynth therefore reports the guard as a distance rule, never as privacy: \texttt{relsynth privacy} (Section~\ref{sec:v25}) remains the test and DP mode (Section~\ref{sec:v25dp}) the guarantee, and the scorecard marks the \texttt{dcr\_ratio} of a guarded table as not independent of the generation. Check C1723 builds the example in one dimension and finds the empty ball.

\setcounter{proposition}{1723}
\begin{proposition}[Shape shrinkage keeps the segment's location]\label{prop:segshape}
In shape mode (\texttt{marginals.segment\_prior\_mode: shape}, the default after 2.2.0) discrete columns keep the weight $a_0$, and a continuous column of a segment with $n\ge1$ observed values, sample median $m_g$ and interquartile range $q_g$ uses, in place of $F_0$, the law $F_0^g$ of $\min(\max(m_g+s_g(X-m_0),\,\ell),\,u)$ with $X\sim F_0$, where $m_0$, $q_0>0$, $\ell$, $u$ are the median, interquartile range, smallest and largest knot of $F_0$, and $s_g=\max(q_g,\,q_0a/(4(n+a)))/q_0$. (a) The mixture and draw-count statements of Props.~\ref{prop:segshrink} and \ref{prop:seglarge} hold with $F_0^g$ in place of $F_0$; the 2.1 output at $a=a_0$ is not inherited, since $F_0^g\ne F_0$. The code therefore uses table mode (the plain $F_0$) whenever \texttt{segment\_prior\_rows} $\le a_0$, which gives the 2.1 output in every mode. (b) If $\ell\le m_g\le u$, the posterior mean predictive $w\,e+(1-w)F_0^g$ has median $m_g$: shrinkage moves no such segment's median in the mean. If $m_g$ lies outside $[\ell,u]$ (a coarse parent, or a segment beyond the table's knots), $F_0^g$ sits at $\ell$ or $u$ and pulls the median toward the range. A posterior draw (Bayesian bootstrap weights on $e$ and on the prior's pseudo-points) varies around it, as without shrinkage, and the fitted marginal matches it up to the interpolation between grid knots. (c) With \texttt{marginals.tail\_trim} at 0 (the default), every non-null synthetic value of the column lies in the smallest interval that contains $[\ell,u]$ and the segment's own values. (A positive \texttt{tail\_trim} extrapolates the trimmed tails linearly on purpose, so as not to expose the extremes, and its values can fall outside.) The parent $F_0$ may be the table's marginal or a coarse group's. (d) No random number is drawn to build $F_0^g$: at a fixed $a$, shape mode draws the same variates from the same generator state as the plain prior, and the random stream is unchanged. (e) Since $F_0^g$ depends on $a$, Prop.~\ref{prop:seglarge} compares two shape-mode predictives only for a fixed $F_0^g$; for any two prior laws $G$, $H$ and $a_0\le a$, the posterior predictive laws satisfy $\mathrm{TV}\big(w_a e+(1-w_a)G,\;w_{a_0}e+(1-w_{a_0})H\big)\le a/(n+a)$. In particular the shape-mode predictive differs from the 2.1 predictive ($a_0$, $F_0$) by at most $a/(n+a)$ in total variation. The bound is for the predictive mixtures; the fitted marginal also interpolates between its grid knots, which the bound does not cover.
\end{proposition}
\begin{proof}
(a) Both proofs use only that the predictive is the mixture $w\,e+(1-w)F_0$ and the draw counts; they hold for any prior law, in particular $F_0^g$. (b) $x\mapsto\min(\max(m_g+s_g(x-m_0),\ell),u)$ is non-decreasing and maps $m_0$ to $m_g$ when $\ell\le m_g\le u$, so $F_0^g$ has median $m_g$; $e$ has median $m_g$ by definition; a mixture of two laws with the same median has that median. (c) The predictive's support is the union of the supports of $e$ and $F_0^g\subseteq[\ell,u]$; the inverse transform interpolates between knots of that union, so it stays in the smallest interval containing both (an outlying segment value can give interpolated values outside $[\ell,u]$ that are not one of the segment's values). Bootstrap weights change the weights of the knots, not their positions. (d) $m_g$, $q_g$, $m_0$, $q_0$ are deterministic functions of the data and of $F_0$'s knots. (e) $w_a\le w_{a_0}$, so the second law is $w_a e+(1-w_a)H'$ with $H'=\big((w_{a_0}-w_a)e+(1-w_{a_0})H\big)/(1-w_a)$, a law. The two laws share the component $w_a e$, so their total variation is $(1-w_a)\,\mathrm{TV}(G,H')\le1-w_a=a/(n+a)$.
\end{proof}
\noindent Shrinking the whole segment toward $F_0$ (\texttt{segment\_prior\_mode: table}, Prop.~\ref{prop:segshrink} as such) also moves its location, and with it what the fact tables read from the segment (which products sell, at what price); on the fake-data suite it fails C09, C19 and C43. Shape mode keeps the location in the posterior mean for segments whose median lies in the parent's knot range (b, up to the interpolation between grid knots) and the range (c), and those checks pass. Check C1724 fits segments far from the table's centre and checks (b) on the fitted posterior mean marginal (\texttt{bootstrap=False}, median within 2\% of the table's range), (c) and (d) at a fixed $a$, and a segment whose median lies beyond $u$ (its predictive median falls below $m_g$), and checks (e) for random discrete $e$, $G$, $H$, $n$, $a_0$ and $a$; a second check builds a segment whose values leave the interval of (c) with a positive \texttt{tail\_trim}.


\subsection*{Data-driven prior strength and the privacy floor (v2.3)}
Props.~\ref{prop:segshrink}--\ref{prop:segshape} take the prior weight $a$ as a fixed number of rows. A fixed number is a tuning constant: 20 rows was chosen in 2.2 on two named databases, and 40 rows hurt the price fidelity of Contoso's fact tables. From 2.3 the weight is a parameter of the hierarchical model, estimated from the segments themselves for each column and each level of the hierarchy (\texttt{marginals.segment\_prior\_rows: auto}, code \texttt{relsynth/shrinkage.py}; off by default and turned on per table in the spec, since its fidelity cost depends on the schema), and privacy enters as a separate, explicit floor on the number of rows of mass in every segment marginal (\texttt{marginals.segment\_floor\_rows}). The estimate is parametric empirical Bayes \cite{efronmorris1975,morris1983}; the floor is a threshold rule in the sense of statistical disclosure control \cite{hundepool2012}, stated on the posterior predictive instead of on released cells.

\setcounter{proposition}{1724}
\begin{proposition}[Empirical-Bayes concentration]\label{prop:segeb}
Let the segments $g=1,\dots,G$ of an entity table (inside reference groups $c=1,\dots,C$: the coarse groups, or one group) have $n_g$ rows, and let the segment laws of a column be independent draws $F_g\sim\mathrm{DP}(a,F_c)$ given the reference law, the rows of a segment i.i.d.\ from $F_g$. (a) For every threshold $t$ (or level $\ell$ of a discrete column) two distinct rows of one segment have $\mathrm{corr}(\mathbf 1\{x_i\le t\},\mathbf 1\{x_j\le t\})=\rho=1/(a+1)$, whatever $t$ and $F_c$ are. (b) For indicators $Y_{gi}$ of one threshold with segment means $\bar Y_g$, reference means $\bar Y_c$, $N=\sum_g n_g$, $\mathrm{SSB}=\sum_g n_g(\bar Y_g-\bar Y_{c(g)})^2$, $\mathrm{SSW}=\sum_g n_g\bar Y_g(1-\bar Y_g)$, $\mathrm{MSB}=\mathrm{SSB}/(G-C)$, $\mathrm{MSW}=\mathrm{SSW}/(N-G)$ and $n_0=(N-\sum_g n_g^2/N_{c(g)})/(G-C)$, the components $\hat\sigma^2_b=(\mathrm{MSB}-\mathrm{MSW})/n_0$ and $\hat\sigma^2_w=\mathrm{MSW}$ are unbiased for $\rho\,\sigma^2$ and $(1-\rho)\,\sigma^2$, with $\sigma^2$ the variance of the indicator, given the reference laws. (c) With $\hat\rho=\sum_t\hat\sigma^2_{b,t}/\sum_t(\hat\sigma^2_{b,t}+\hat\sigma^2_{w,t})$ over any set of thresholds (or the levels), $\hat a=(1-\hat\rho)/\hat\rho$ is consistent for $a$ as $G\to\infty$ with bounded $n_g\ge2$ when the reference laws agree at those thresholds ($F_c(t)$ the same for every $c$: one reference group, or coarse groups that differ elsewhere); with reference laws that differ there, the pooled mean squares weight the groups' indicator variances differently and $\hat a$ is a pooled estimate that need not converge to $a$. (d) The weight $n_g/(n_g+a)$ that the posterior predictive gives the segment's own rows is the Bühlmann credibility factor with $k=\sigma^2_w/\sigma^2_b=a$: among estimators of $F_g(t)$ linear in the segment's empirical distribution function and $F_c(t)$, it minimises the mean squared error, for every $t$ at once. (e) Computing $\hat a$ draws no random number; so at the same weights the fit is the fixed-weight fit of Props.~\ref{prop:segshrink}--\ref{prop:segshape}, and every statement there holds with $a=\hat a$ (the per-column weight replaces the common one).
\end{proposition}
\begin{proof}
(a) By Ferguson's theorem \cite{ferguson} $F_g(A)\sim\mathrm{Beta}(aF_c(A),a(1-F_c(A)))$ for $A=(-\infty,t]$ (or $\{\ell\}$), so with $p=F_c(A)$: $E[F_g(A)]=p$ and $\mathrm{Var}F_g(A)=p(1-p)/(a+1)$. For $i\ne j$ in segment $g$, $P(x_i\in A,x_j\in A)=E[F_g(A)^2]=p^2+p(1-p)/(a+1)$, so the covariance of the indicators is $p(1-p)/(a+1)$ and their variance is $p(1-p)$; the correlation is $1/(a+1)$. For a discrete column the vector of level probabilities is $\mathrm{Dir}(aF_c)$ and each level indicator is the Dirichlet-multinomial case \cite{mosimann1962} with the same $\rho$. (b) The indicators follow the nested random-effects model $Y_{gi}=\mu_c+b_g+e_{gi}$ with $E b_g=0$, $\mathrm{Var}\,b_g=\rho\sigma^2$, $\mathrm{Var}\,e_{gi}=(1-\rho)\sigma^2$, uncorrelated, which by (a) is exactly the covariance of the rows. The expected mean squares of the unbalanced nested one-way analysis of variance are $E\,\mathrm{MSW}=\sigma^2_w$ and $E\,\mathrm{MSB}=\sigma^2_w+n_0\sigma^2_b$ (Sokal and Rohlf \cite{sokalrohlf}, Sect.~10.3; within groups $\sum_i(Y_{gi}-\bar Y_g)^2=n_g\bar Y_g(1-\bar Y_g)$ because $Y^2=Y$), which gives both unbiasedness statements. (c) Numerator and denominator are sums of segment-level terms with bounded moments; by the law of large numbers over segments each converges to its expectation, $\rho\sum_t\sigma^2_t$ and $\sum_t\sigma^2_t$ (one $\sigma^2_t=F_c(t)(1-F_c(t))$ for every reference group under the hypothesis of (c)), so $\hat\rho\to\rho$ and $\hat a\to a$ by continuity of $\rho\mapsto(1-\rho)/\rho$ on $(0,1]$. Summing over thresholds before taking the ratio is the estimator of Weir and Cockerham \cite{weircockerham1984} across alleles. (d) Bühlmann's credibility theorem \cite{buhlmann1967,buhlmannstraub1970} for the structure (a): the best linear estimator of $E[Y\mid g]=F_g(A)$ from $\bar Y_g$ and $p$ is $Z\bar Y_g+(1-Z)p$ with $Z=n_g/(n_g+\sigma^2_w/\sigma^2_b)$, and $\sigma^2_w/\sigma^2_b=(1-\rho)/\rho=a$ for every $A$. It equals the posterior mean $(n_g\bar Y_g+ap)/(n_g+a)$ of the DP. (e) $\hat a$ is a deterministic function of the codes and of the segment assignment.
\end{proof}
\noindent The estimate is truncated to $[0,N]$: $\hat\rho\le0$ (segments no more alike within than between) gives full pooling at $N$ rows, the prior's own limit; $\hat\rho\ge1$ gives $a_0$. In shape mode the rows are first centred at their segment's median and scaled by its interquartile range (segments with at least four values and a positive range), so $\hat a$ measures how far the segments' \emph{shapes} differ, which is what shape mode shrinks; a segment's location is not evidence against pooling its shape. The estimate is a plug-in: the uncertainty of $\hat a$ is not propagated (a full hierarchical treatment would put a prior on $a$, Escobar and West \cite{escobarwest1995}). Check C1725 simulates segments from DP hierarchies with known $a$ and recovers it, checks the balanced-design formula, the nested reference and the absence of random draws.

\setcounter{proposition}{1725}
\begin{proposition}[Privacy floor]\label{prop:segfloor}
With \texttt{segment\_floor\_rows} $=k$, a segment of $n$ rows (non-null values for a continuous column) uses the weight $a_g=\max(a_0,\hat a,k-n)$ on every column that \texttt{segment\_prior\_mode} covers (in shape mode the continuous columns; discrete columns keep $a_0$ and their own shares, so the bound below is per covered column). (a) In the posterior predictive each source row of the segment has weight $1/(n+a_g)\le1/k$ in the segment's own part, and in a posterior draw that weight has mean $1/(n+a_g)$; the parent law $F_c$ the segment is pooled with is fitted on all $N_c$ rows of the table or coarse group, the segment's own included, and adds at most $a_g/((n+a_g)N_c)\le1/N_c$ for a row, so a row carries at most $1/k+1/N_c$ in all. (b) For a segment of $n\ge k$ rows the floor is inactive and the fit is that of Prop.~\ref{prop:segeb}. (c) Changing one source row of a segment changes the posterior predictive by at most $1/(n+a_g)\le 1/k$ in total variation, when the change moves neither the parent law nor $\hat a$ (a fixed \texttt{segment\_prior\_rows}) nor, in shape mode, the segment's median and interquartile range that place the prior (a fixed transformed prior); and the parent law $F_c$, itself a mixture over at least $N_c$ rows, moves by at most $1/N_c$. (d) The floor is not differential privacy: (c) bounds the predictive, not the probability of an output, and the segment assignment and $\hat a$ themselves depend on every row.
\end{proposition}
\begin{proof}
(a) The posterior predictive is $\frac{1}{n+a_g}\sum_{i}\delta_{x_i}+\frac{a_g}{n+a_g}F_c$ and $n+a_g\ge k$, with $n$ the observations the fit uses (non-null values, at most \texttt{max\_obs} of them, for a continuous column). In a draw the segment's rows share $V\sim\mathrm{Beta}(n,a_g)$ with Dirichlet$(1,\dots,1)$ weights, so a row's weight has mean $E[V]/n=1/(n+a_g)$. A row has weight at most $1/N_c$ in $F_c$, which enters with weight $a_g/(n+a_g)\le1$. (In shape mode $F_c$ enters moved and scaled, so its part is not an atom at the row's value; the bound is that of its mass.) (b) $k-n\le0$. (c) Replacing $x_i$ by $x_i'$ changes the mixture by $\frac{1}{n+a_g}(\delta_{x_i'}-\delta_{x_i})$, of total variation at most $1/(n+a_g)$; the same argument applies to $F_c$ with its $N_c$ rows. (d) A bounded change of a predictive law is compatible with outputs whose probability ratio under neighbouring tables is unbounded (the inverse transform maps a point mass to a point); Prop.~\ref{prop:ncnotdp} gives the general reason. DP mode (Section~\ref{sec:v25dp}) is the guarantee.
\end{proof}
\noindent The floor is the explicit privacy control of small segments; it is a policy choice of the deploying organisation, like $k$ in $k$-anonymity \cite{sweeney2002} or the minimum cell frequency of a threshold rule \cite{hundepool2012}, and it is never inferred from the data. The empirical-Bayes weight serves fidelity: where segments really differ (a product's price fixes which facts it sells to), $\hat a$ is small and the segments keep their own laws; where they do not, they are pooled. Check C1726 fits segments of 3 to 60 rows and checks the bound of (a) on the fitted posterior mean marginal, (b) byte for byte, and (c) by replacing one row at a fixed weight and a fixed prior.

\setcounter{proposition}{\value{savedpropgthree}}
